\chapter*{Resumen}
\addcontentsline{toc}{chapter}{Resumen}

A principios del año 2024 se abrió una nueva línea de investigación dentro del Grupo de Física de la Atmósfera, orientada a indagar la posible existencia de vida microbiana en venas líquidas del hielo \cite{mader2006}, \cite{rohde2007}, \cite{dani2012}. En sus primeros estadios, la investigación se ha centrado en el estudio del movimiento de microorganismos presentes en aguas estancadas de la ciudad de Córdoba, en el rango de temperaturas [0 °C; +15 °C], con el objetivo de comprender, entre otros aspectos, su posible supervivencia en entornos cercanos al punto de fusión del hielo \cite{pedernera2024}.

En particular, se ha puesto el foco en los tipos de trayectorias que describen estos microorganismos y en sus velocidades medias, así como en la dependencia de estas magnitudes con la temperatura del medio. El seguimiento se realizó mediante filmaciones obtenidas con una cámara acoplada a un microscopio óptico utilizado para visualizar las muestras. En todas las muestras analizadas se identificaron cuatro tipos de microorganismos.

Actualmente, el análisis de los movimientos se realiza de forma completamente manual mediante un software de uso gratuito (Tracker), lo que implica una elevada inversión de tiempo y limita el avance de la investigación. En este contexto, el desarrollo de un sistema automático de detección y seguimiento resulta fundamental para optimizar el procesamiento de los datos experimentales y permitir la obtención de métricas cuantitativas adicionales, tales como distribuciones de velocidad, trayectorias individuales, variaciones en la aceleración y cambios en los patrones de movimiento frente a estímulos externos (por ejemplo, modificaciones en la temperatura del medio).
